\documentclass[12pt]{article}
\usepackage[a4paper,margin=1.25cm]{geometry}
\usepackage{amsmath,amssymb,mathtools}
\usepackage{graphicx}
\usepackage{booktabs}
\usepackage{tabularx}
\usepackage{tikz}
\usepackage[T1]{fontenc}
\usepackage[utf8]{inputenc}
\usepackage[english]{babel}

\setlength{\parindent}{0pt}
\setlength{\parskip}{0.45em}
\setlength{\abovedisplayskip}{5pt plus 1pt minus 1pt}
\setlength{\belowdisplayskip}{5pt plus 1pt minus 1pt}
\setlength{\abovedisplayshortskip}{3pt plus 1pt minus 1pt}
\setlength{\belowdisplayshortskip}{3pt plus 1pt minus 1pt}

\newcommand{\dd}{\,\mathrm{d}}
\newcommand{\e}{\mathrm{e}}
\newcommand{\Sn}{S_n}
\newcommand{\Tn}{T_n}
\newcommand{\In}{I_n}
\newcommand{\D}{\Delta}
\newcommand{\Hh}{H}

\begin{document}

\begin{center}
{\Large\bfseries An alternating binomial sum of logarithms}\\[0.35em]
{\large a high finite difference is a tiny integral}
\end{center}

\[
\boxed{\displaystyle
\Sn=\sum_{k=1}^{n-1}(-1)^k\binom nk\log(k+1)
=(-1)^{n+1}\log(n+1)-\In,
\qquad
\In=\int_0^1\frac{(1-x)^n}{-\log x}\dd x,}
\]
\[
\boxed{\displaystyle
\In=\frac1{n\log n}\left(1-\frac{\gamma}{\log n}
+O\!\left(\frac1{\log^2 n}\right)\right),
\qquad 0<\In<\frac1n .}
\]

\textbf{Overview.}
The alternating binomial coefficients are the \(n\)-th finite difference
operator in disguise:
\[
\D^{n}a_0=\sum_{j=0}^{n}(-1)^{n-j}\binom nj a_j .
\]
Applying this to \(a_m=\log(m+1)\), the \emph{completed} sum
\(\Tn=\sum_{k=0}^{n}(-1)^k\binom nk\log(k+1)\) is exactly
\((-1)^n\D^n a_0\). High-order differences of a slowly growing smooth
function are small, essentially because \(\log\) is nearly a
polynomial of low degree on any short scale. Here the smallness is
quantified exactly: \(\Tn=-\In\) with \(\In\) an explicit positive
integral of size \(\frac1{n\log n}\).

The sum \(\Sn\) of the problem omits the terms \(k=0\) (which is
\(\log1=0\)) and \(k=n\) (which is \((-1)^n\log(n+1)\)), so
\[
\Sn=\Tn-(-1)^n\log(n+1)=(-1)^{n+1}\log(n+1)-\In .
\]
Thus \(\Sn\) is dominated by the single omitted term \(\pm\log(n+1)\),
corrected by a quantity that vanishes like \(\frac1{n\log n}\).

Two independent derivations of \(\In\) are given: a discrete one, through
the finite-difference calculus of \(\frac1{m+\alpha}\), and a continuous
one, through the Frullani-type representation of \(\log(k+1)\). The first
gives the sharp elementary bound \(\In<\frac1n\); the second explains why
Euler's constant appears in the second-order term.

\section*{1. The sum is a finite difference}

Let \(a_m=\log(m+1)\) and \(\D a_m=a_{m+1}-a_m\). Iterating and using
Pascal's rule \(\binom{r+1}{j}=\binom rj+\binom r{j-1}\),
\[
\D^{r}a_m=\sum_{j=0}^{r}(-1)^{r-j}\binom rj a_{m+j}.
\tag{1}
\]
Setting \(m=0\), \(r=n\) and multiplying by \((-1)^n\),
\[
\Tn:=\sum_{k=0}^{n}(-1)^k\binom nk\log(k+1)=(-1)^n\D^{n}a_0 .
\tag{2}
\]
The sum of the problem omits the two extreme terms:
\[
k=0:\ \log1=0,
\qquad
k=n:\ (-1)^n\log(n+1),
\]
so
\[
\boxed{\displaystyle \Sn=\Tn-(-1)^n\log(n+1).}
\tag{3}
\]
Everything therefore reduces to showing that \(\Tn\) is small, i.e. that
the \(n\)-th finite difference of \(\log(m+1)\) at \(m=0\) is small.

\section*{2. Discrete evaluation: the harmonic ruler}

Write the first difference as an average:
\[
\D a_m=\log\left(1+\frac1{m+1}\right)=\int_0^1\frac{\dd t}{m+1+t}.
\tag{4}
\]
This is the key step: it expresses \(\D a_m\) as a superposition of the
elementary sequences \(m\mapsto\frac1{m+\alpha}\), for which the finite
differences are known exactly:
\[
(-1)^{r}\D^{r}\frac1{m+\alpha}\bigg|_{m=0}
=\frac{r!}{\alpha(\alpha+1)\cdots(\alpha+r)},
\qquad\alpha>0 .
\tag{5}
\]
(Immediate by induction, since
\(\frac1{m+\alpha}-\frac1{m+1+\alpha}=\frac1{(m+\alpha)(m+1+\alpha)}\).)

Applying \(\D^{n-1}\) to (4) and using (5) with \(\alpha=1+t\),
\(r=n-1\),
\[
\In:=(-1)^{n-1}\D^{n}a_0
=\int_0^1\frac{(n-1)!\dd t}{(1+t)(2+t)\cdots(n+t)} .
\tag{6}
\]
By (2), \(\Tn=-\In\), which with (3) gives the first boxed formula.

Since the integrand of (6) is positive and, for \(t\in[0,1]\), bounded by
its value at \(t=0\),
\[
\boxed{\displaystyle
0<\In<\frac{(n-1)!}{1\cdot2\cdots n}=\frac1n .}
\tag{7}
\]

\textbf{The harmonic ruler.} Bound (7) is best understood by comparison
with the harmonic numbers. Let \(\Hh_0=0\) and
\(\Hh_m=1+\frac12+\cdots+\frac1m\), so that
\(\D \Hh_m=\frac1{m+1}\). By (5) with \(\alpha=1\),
\[
(-1)^{n-1}\D^{n}\Hh_0=\frac1n .
\tag{8}
\]
Comparing (4) with \(\D \Hh_m=\frac1{m+1}\): the logarithm's first
difference is the \emph{average} of \(\frac1{m+1+t}\) over
\(t\in[0,1]\), hence strictly smaller than \(\frac1{m+1}\). Since all the
relevant differences carry the same alternating sign, the same ordering
persists at order \(n\), which is exactly (7). In words: \(\log(m+1)\)
and \(\Hh_m\) have the same alternating-difference structure, and the
logarithm sits below the harmonic ruler.

\section*{3. Continuous evaluation: a Frullani representation}

An independent route gives a form better suited to asymptotics. For
integers \(k\ge0\),
\[
\log(k+1)=\int_0^1\frac{1-x^k}{-\log x}\dd x,
\tag{9}
\]
which is the Frullani integral \(\int_0^1\frac{x^{b}-x^{a}}{\log x}\dd x
=\log\frac{b+1}{a+1}\) at \(a=0\), \(b=k\).

Insert (9) into \(\Tn\) and exchange the finite sum with the integral.
Since \(\sum_{k=0}^{n}(-1)^k\binom nk=0\) for \(n\ge1\) and
\(\sum_{k=0}^{n}(-1)^k\binom nk x^k=(1-x)^n\),
\[
\Tn=\int_0^1\frac{0-(1-x)^n}{-\log x}\dd x
=-\int_0^1\frac{(1-x)^n}{-\log x}\dd x .
\]
Hence
\[
\boxed{\displaystyle
\In=\int_0^1\frac{(1-x)^n}{-\log x}\dd x
=\int_0^1\frac{(n-1)!\dd t}{(1+t)(2+t)\cdots(n+t)} ,}
\tag{10}
\]
the two representations agreeing, as one checks numerically to full
precision for any \(n\).

\section*{4. Asymptotics: where \(\gamma\) comes from}

Use the integral representation. Substituting \(x=\frac yn\) and writing
\(L=\log n\),
\[
\In=\frac1n\int_0^n\frac{(1-y/n)^n}{L-\log y}\dd y .
\tag{11}
\]
Two expansions now govern the answer:
\[
(1-y/n)^n\longrightarrow\e^{-y},
\qquad
\frac1{L-\log y}=\frac1L\left(1+\frac{\log y}{L}
+O\!\left(\frac{\log^2 y}{L^2}\right)\right).
\tag{12}
\]
Therefore
\[
\In=\frac1{nL}\left(
\int_0^\infty \e^{-y}\dd y
+\frac1L\int_0^\infty \e^{-y}\log y\dd y
+O\!\left(\frac1{L^2}\right)\right).
\]
The two integrals are classical:
\[
\int_0^\infty \e^{-y}\dd y=1,
\qquad
\int_0^\infty \e^{-y}\log y\dd y=\Gamma'(1)=-\gamma .
\tag{13}
\]
Hence
\[
\boxed{\displaystyle
\In=\frac1{n\log n}\left(1-\frac{\gamma}{\log n}
+O\!\left(\frac1{\log^2 n}\right)\right).}
\tag{14}
\]

\textbf{The same constant, discretely.} Formula (6) reproduces (14)
without any Gamma function. Factoring out \((n-1)!\),
\[
\In=\frac1n\int_0^1\prod_{r=1}^{n}\left(1+\frac tr\right)^{-1}\dd t
=\frac1n\int_0^1\exp\left(-\sum_{r=1}^{n}\log\left(1+\frac tr\right)\right)\dd t,
\tag{15}
\]
and for \(0\le t\le1\),
\[
\sum_{r=1}^{n}\log\left(1+\frac tr\right)=t\,\Hh_n+O(t^2).
\]
The integrand is therefore concentrated on \(t\asymp \Hh_n^{-1}\);
substituting \(t=u/\Hh_n\),
\[
\In=\frac1{n\Hh_n}\int_0^{\Hh_n}\e^{-u}
\exp\!\left(O\!\left(\frac{u^2}{\Hh_n^2}\right)\right)\dd u
=\frac1{n\Hh_n}\left(1+O\!\left(\frac1{\Hh_n^{2}}\right)\right).
\tag{16}
\]
Expanding \(\Hh_n=\log n+\gamma+O(n^{-1})\) turns (16) into (14). So the
\(\gamma\) in the second-order term is, on the discrete side, nothing but
\(\Hh_n-\log n\to\gamma\), and on the continuous side the value
\(\Gamma'(1)\). The two occurrences are the same constant seen from the
two ends of (10).

\section*{5. Numerical check}

\begin{center}
\small
\renewcommand{\arraystretch}{1.35}
\begin{tabularx}{0.95\textwidth}{@{}lXXXX@{}}
\toprule
\(n\) & \(\In\) (exact) & \(\frac1{n\log n}\bigl(1-\frac\gamma{\log n}\bigr)\) & \(\frac1{n\Hh_n}\) & bound \(\frac1n\) \\
\midrule
\(5\)    & \(0.0871265\)   & \(0.0796993\)   & \(0.0875912\)   & \(0.2\) \\
\(20\)   & \(0.0145850\)   & \(0.0134745\)   & \(0.0138976\)   & \(0.05\) \\
\(100\)  & \(0.00201125\)  & \(0.00189930\)  & \(0.00192776\)  & \(0.01\) \\
\(1000\) & \(0.000137012\) & \(0.000132668\) & \(0.000133592\) & \(0.001\) \\
\bottomrule
\end{tabularx}
\end{center}

The discrete approximation \(\frac1{n\Hh_n}\) of (16) is visibly better
than the two-term logarithmic expansion (14), which is expected: the
expansion parameter of (14) is \(\frac1{\log n}\), still \(0.14\) at
\(n=1000\), whereas (16) resums the whole harmonic number.

\section*{6. Summary of the identity}

Collecting (3), (10) and (14),
\[
\Sn=\sum_{k=1}^{n-1}(-1)^k\binom nk\log(k+1)
=(-1)^{n+1}\log(n+1)
-\frac1{n\log n}\left(1-\frac{\gamma}{\log n}
+O\!\left(\frac1{\log^2 n}\right)\right).
\]


\textbf{The sum is not small; the completed sum is.} \(\Sn\) grows like
\(\log n\) in absolute value and alternates in sign with \(n\), whereas
\(\Tn=\Sn+(-1)^n\log(n+1)\) tends to \(0\). Completing the sum by the
term \(k=n\) is what turns it into a finite difference; the binomial
coefficients only form a difference operator when the full range is
present.

\textbf{The cancellation is enormous.} The individual terms of \(\Tn\)
reach size \(\binom{n}{n/2}\log n\approx2^n\), while the total is
\(\frac1{n\log n}\). Formula (10) explains the mechanism: the alternating
sum of \(x^k\) with binomial weights is \((1-x)^n\), which is
exponentially small everywhere except near \(x=0\), where the factor
\(\frac1{-\log x}\) is itself small. The competition between these two
factors, resolved in Section 4, produces the scale
\(\frac1{n\log n}\): an algebraic \(\frac1n\) from the width of the
region where \((1-x)^n\) survives, and a logarithmic \(\frac1{\log n}\)
from the weight there.

\end{document}
